\cotefile{156-9}
\scanpage{34}
\mspage{30}
% archive sheet number: 33
% margin note written sideways
\margin{\uline{8 juillet}}
% source lines 1--3: heading
\uline{Digression 3} Magmas \unsure{munis} de \unsure{strates} \inserted{dites} \uline{compactes} (\uline{Magmas \unsure{compactologisés|compactologiques}}).

\uline{Notion de figure localement finie}

% source lines 4--7
% slanted
\margin{\unsure{La donnée} de $M_c$ s'appelle une ``\uline{compactologie}''\unsure{.}}
On se donne une partie $M_c$ de $M$, (structure supplémentaire) formée des \uline{strates} dites \deleted{``admissibles''} ``\uline{compactes}''.

\deleted{Ax} L'axiome \unsure{préalable} est le suivant

% source lines 8--10
\uline{Mag cpt 1}\quad \deleted{a) Si $X\in M$, \unsure{et} $\Rightarrow$ alors $\widetilde X$ fini.}
\[
\begin{aligned}
\text{a)}\quad &X\in M_c\ \Longrightarrow\ \widetilde X\ \text{fini.}\\
\text{b)}\quad &X\in M_c,\ Y\ll X\ \Longrightarrow\ Y\in M_c
\end{aligned}
\]
% source lines 10--13
--- en d'autres termes : $M_c$ est fermé dans $(M,\ll)$, ce qui revient au même, dans $(M,\overset{\circ}{\ll})$ et dans $(M,\trianglelefteq)$.

% slanted; left margin opposite lines 11--19
\margin{Un \unsure{magma|magasin} \unsure{compactologisé|compactologique} est dit ``à strates compactes'' si $M_c=M$. Il est dit \uline{compact} si $\exists F\in\unsure{\mathcal F|\mathscr F|\mathfrak F}(M)$, \unsure{à un} $\widetilde F$ \uline{fini}, tel que $\operatorname{supp}^{\circ}F=\unsure{M_c|M}$\unsure{.} i.e. $\nexists X\in M,\ F\mathrel{\mid\circ\mid}\deleted{\unsure{X}}X$.}
% slanted; left margin below the preceding note
\margin{\uline{Figure \unsure{compacte}} : $\widetilde F$ fini, et $\subset M_c$. \uline{\unsure{Préfigures|Préfigure}} ``\uline{compactes}'' : \deleted{\unsure{finies d'ordres}} une figure compacte.}

% source lines 14--17
Par la suite, on sera amené \unsure{sûrement} d'introduire un axiome Mag cpt 2, du type Borel-Lebesgue ; \inserted{et aussi que $M_c$ est ``riche'' dans $M$ (\unsure{p.v.|p.r.} aux subdivisions)}

% source lines 18--20
Quand on se donne de plus un \unsure{atelier} $\mathcal F$, \unsure{déterminant} un magasin, on exige

\deleted{Mag cpt (at)}

% source lines 21--24
\uline{Mag at cpt}\quad Si $F\ll X$, avec $F\in\mathcal F$, $X\in M_c$, alors $\widetilde F$ fini. Mieux : Si $X\in M_c$, alors $\forall F\in\mathcal F$, l'ensemble des $\deleted{X}Y\in\widetilde F$ telles que $\deleted{\widetilde F\mathrel{\mid\circ\mid}X}\ Y\parallel X$ est fini,
% left margin; continuation of the preceding sentence
\margin{i.e. $\widetilde F\mathrel{\mid\circ\mid}\widetilde X$, a un complémentaire fini dans $\widetilde F$.}

% source lines 25--28: paragraph cancelled with diagonal strokes
\deleted{On note ceci : Pour $M_c$ donné, \inserted{une compactologie donnée} il existe un $\mathcal F$ \unsure{attenant} à $M$, le plus grand possible \unsure{formé} \unsure{aux|avec|des} satisfaisant la condition envisagée. Les figures que}
\scanpage{35}
\mspage{31}
% archive sheet number: 34
% source lines 1--4
Pour $M_c$ donné, et une partie $\Phi\subset M$, on dira que $\Phi$ est \uline{localement finie} si \deleted{$\forall X\in M_c,$}
\[
\forall X\in M_c,\quad
\text{\deleted{posant}\inserted{l'ens.}}\quad
\Phi_X=\{Y\in\Phi\mid Y\mathrel{\overline{\parallel}}X\}
\quad\text{\inserted{est fini} \deleted{on a $\Phi\unsure{\setminus|,}\Phi_X$ fini.}}
\]
% source line 5, closely written above the next paragraph
Les parties loc finies de $M$ forment un \uline{antifiltre}.

% source lines 6--9
L'ensemble \deleted{des} $\unsure{\mathcal F|\mathscr F|\mathfrak F}_{l.f}(M)$ des figures loc. finies de $M$ forme un \inserted{$M$-}atelier, à condition pourtant de supposer \deleted{que} ceci (qui assure que $M\subset\mathcal F_{l.f}(M)$)

% source lines 10--12
\uline{Mag cpt 1}\quad c) \deleted{$\forall\unsure{X|Y}\in M_c,$} $X\in M$, $\widetilde X$ est loc fin pour $M_c$ i.e. $\forall Y\in M_c$, l'ensemble \deleted{des}
\[
\{X'\in\widetilde X\mid X'\mathrel{\overline{\parallel}}Y\}\quad\text{est fini.}
\]

% source lines 13--14
Si $\Phi$ est une préfigure loc finie, alors $\overline\Phi^{(\trianglelefteq)}$ est encore loc finie.

% source lines 15--18
L'axiome Mag at cpt signifie que $\mathcal F\subset\mathcal F_{\mathrm{loc\,fin.}}$. Ainsi $\mathcal F_{\mathrm{loc\,fin}}$ est le plus grand atelier \unsure{adhérant|dérivant|attenant} à $M$, qui soit compatible à la compactologie.

% source lines 19--24
Inversement, donnons nous $(M,\trianglelefteq,\overset{\circ}{\ll},\mathrel{\mid\circ\mid},\mathcal F)$. On va construire une \deleted{une} compactologie la plus grande possible, pour laquelle $M_c,\mathcal F$ soient compatibles, en posant
\[
\begin{aligned}
M_c=\{X\in M\mid {}&\deleted{\unsure{\exists}}\forall F\in\mathcal F,\ \text{l'ens }\{Y\in\widetilde F\mid X\mathrel{\overline{\parallel}}Y\}\\
&\text{est fini}\}.
\end{aligned}
\]
\scanpage{36}
\mspage{32}
% archive sheet number: 35
% source lines 1--3
Il est clair que pour cet $M_c$, les conditions Mag cpt 1 a) b) c) sont satisfaites (pour a)\inserted{ et c)}, on prendra $F=\widetilde X$).

% source lines 4--10
\deleted{\unsure{A}} Si $\unsure{\mathcal F|\mathscr F|\mathfrak F}$ est formé des figures finies, ledit $M_c$ est \unsure{ainsi|alors} la plus grande des compactologies sur $M$. Mais il ne semble pas raisonnable de \uline{définir} en général, $M_c$ en termes du seul \unsure{magasin} $(M,\trianglelefteq,\overset{\circ}{\ll},\mathrel{\mid\circ\mid})$, par les conditions précédentes, comme \deleted{\unsure{les}}
\[
M_c=\{X\in M\mid\forall Y\in M,\ \text{l'ens }\{Y'\in\widetilde Y\mid Y'\mathrel{\overline{\parallel}}X\}\ \text{fini}\},
\]
% source lines 12--21
laquelle donnerait $M_c=M$ quand toutes les multistrates \deleted{de} $\widetilde Y$ de $M$ sont finies. P. ex. dans le cas d'un ensemble ordonné $I$, muni d'une seule extrémité ou \unsure{sans} extrémité (\uline{intervalle ordonné} qui n'est pas un segment), il est naturel de considérer comme compactes les seules strates de la forme $a\ (a\in I)\ I_{a,b}\ (a,b\in I)$.

% source lines 22--26
Il semble bien que ce soit la donnée supplémentaire de la compactologie, qui soit une structure souple plus intéressante que celle d'un atelier \unsure{attenant} particulier.
\scanpage{37}
\mspage{33}
% archive sheet number: 36
% source lines 1--10
Elle permet, en plus des \deleted{l'atelier des} figures compactes (et des figures finies dont on disposait déjà avant), d'introduire les figures localement finies. Il \unsure{ne|me} semble pas qu'en pratique, il y ait lieu de travailler avec d'autres ateliers attenants que celui des figures finies, celui des figures \uline{localement finies}, ou (à la rigueur) celui de \uline{toutes} les $M$-figures.

% source lines 11--14
\uline{NB} Les \uline{figures compactes} forment un atelier attenant à $M_c$, qui est un sous-\uline{\unsure{atelier|magasin}} de $M$, mais non à $M$, sauf si $M_c\unsure{=|\equiv}M$.

% source lines 15--19
Ainsi, on peut dire que la base \inserted{technique} de l'approche combinatoire de la topologie que je propose, est plutôt \deleted{la} notion de \uline{magasin compactologique}
\[
(M,\trianglelefteq,\overset{\circ}{\ll},\mathrel{\mid\circ\mid},M_c).
\]
% source lines 20--22
C'est une donnée plus simple que celle d'un atelier attenant à $M$, car $M_c\in\unsure{\mathcal P|\mathscr P}(M)$, et $\unsure{\mathcal F|\mathscr F|\mathfrak F}\in\mathcal P(\mathcal P(M))$.
\scanpage{38}
\mspage{34}
% archive sheet number: 37
% upper margin note, enclosed by a long curved bracket
\margin{Terminologie : atelier, magasin, strate, multistrates.}
% source lines 1--7
\uline{Digression 4}\quad Ainsi, de plus en plus, l'accent principal dans le travail de fondements, se porte sur $M$, plutôt que sur $\unsure{\mathcal F|\mathscr F|\mathfrak F}$. Du coup, cela suggère aussi un changement de terminologie. Le \deleted{terme}\inserted{nom} \unsure{initial} principal, d'\uline{atelier}, \uline{devrait être donné à} $M$.
% source lines 7--18
Cet atelier travaille avec les objets qui figurent dans $M$, les \inserted{strates \unsure{ou}} multistrates, pour constituer avec eux des \uline{figures}. L'ensemble de ces figures peut-être considéré comme constituant un \uline{magasin attenant à l'atelier} (plutôt que l'inverse) --- plusieurs magasins peuvent être attenants au même atelier, comme on vient de le rappeler.
% source lines 19--23
Dans le cas d'un magasin \unsure{compactologisé}, on pourrait appeler peut-être $M_c$, formé des (multi)strates compactes, le \uline{coeur} de l'atelier.
\scanpage{39}
\mspage{35}
% archive sheet number: 38
% source lines 1--8
Il reste la \unsure{perplexité}, sur la terminologie parallèle strate -- multistrates. Finalement, je vais appeler les $X\in M$ les \uline{strates} de l'atelier $M$. Chaque strate détermine l'ensemble $\widetilde X$ de ses \uline{sous-strates}, ou \uline{strates incidentes}, C'est $\widetilde X$ qui sera la \uline{multistrate associée à \unsure{les|la}} \uline{strate} $X$, ou déterminée, définie par $X$, \deleted{\unsure{ou la}}
% source lines 9--16
Ainsi, pour une figure $F$, on distingue entre les \uline{strates de} $F$, qui sont les éléments de $\widetilde F$, ou \uline{strates} de $M$ \uline{incidentes} à $F$, et les \uline{multistrates de} $F$, qui sont les $\widetilde X$ pour $X\in\widetilde F$. On pourrait désigner l'ens de ces multistrates par $\widetilde{\widetilde F}$.

% source lines 17--19
Quand on décrit l'atelier des figures \unsure{élémentaires} ensemblistes sur un ensemble $\unsure{L|\mathcal L}$
\[
\text{Fig \unsure{él|el}}(L)=M
\]
% source lines 20--26
il est entendu que ses objets apparaissent non comme des parties, mais \inserted{disons} comme \deleted{des} ensembles de parties de $L$, \deleted{soit $\widetilde X$}. Cela n'empêche qu'un tel ensemble de parties de $L$, considéré comme objet $\in M$, sera appelé ``strate'', et non ``multistrate''. Ainsi,
\scanpage{40}
\mspage{36}
% archive sheet number: 39
% source lines 1--7; continuation from scan 39
alors qu'un objet \inserted{$X$} de $M$ est un \deleted{\unsure{obj}} élément (de nature particulière) de $\unsure{\mathcal P|\mathscr P}(\mathcal P(\unsure{L|\mathcal L}))$, une multistrate $\widetilde X$ de $M$ doit être \unsure{considéré|considérée}, en bonne logique, comme un élément de $\mathcal P(\mathcal P(\mathcal P(L))$ --- un ensemble d'ensembles de parties de $L$.

% source lines 8--23
De plus en plus, $M$ apparaissant comme l'ensemble de base privilégié, les \uline{figures} $F$ apparaissent comme des parties de $M$. Dans ce point de vue, il n'y a donc pas lieu d'introduire un $\widetilde F$, qui ici n'est autre que $F$. Donc il n'y a pas lieu de parler du ``déploiement'' d'une figure. C'est à l'usage qu'on verra s'il y a lieu, par la suite, d'autonomiser $\unsure{\mathcal F|\mathscr F|\mathfrak F}$ \unsure{p.v.|p.r.} à $M$, par l'introduction d'une structure mixte sur $(M,\mathcal F)$, avec une relation d'incidence $\trianglelefteq$ entre $M$ et $\mathcal F$ --- auquel cas la notation $\widetilde F$ reprend un sens et ses droits.
\scanpage{41}
\mspage{37}
% archive sheet number: 40
% source lines 1--4
\unsure{Dês|Dès} maintenant, je dirai ``atelier'' là où je disais ``magasin'', et inversement. Les axiomes Mag 1 \unsure{:|à} Mag 4 deviennent désormais At 1 : At 4, etc.

% source lines 5--10; large blank space precedes this paragraph
Je reviens à la notion des supports (``\unsure{externes|internes}'') associés à un atelier --- je ne vais plus le noter $M$, il faut avoir le courage de ses \unsure{options} mais $\unsure{A|\mathcal A}$, donc
\[
(A,\trianglelefteq,\overset{\circ}{\ll},\mathrel{\mid\circ\mid}).
\]
% source lines 11--18
La donnée de $\mathrel{\mid\circ\mid}$, relation symétrique \unsure{antireflexive|antiréflexive}, nous permet de définir un ensemble de parties particulières de $A$, \unsure{appelées} ``\unsure{supports|support}'', parties fermées pour $\overset{\circ}{\ll}$ (mais non pour $\trianglelefteq$ en général), \deleted{La donnée}
\[
\Sigma_A\subset\unsure{\mathcal P|\mathscr P}(A),
\]
stable par intersections quelconques, qui
